In this study, a methodology was developed to automatically classify ISUP grades in prostate histopathological images, addressing relevant challenges such as class imbalance, label noise, and the ordinal nature of tumor progression. The introduction of a hybrid loss function, combining ordinal regression and focal loss, improved model performance by explicitly capturing the hierarchical relationship among malignancy grades. Notably, the largest performance gain was achieved through the application of the ordinal loss, which, when combined with noise removal, increased the $\kappa_{\text{quad}}$ index by approximately 2.5\%, while maintaining high stability (standard deviation of 0.009). Overall, the proposed approach improved the $\kappa_{\text{quad}}$ index from 0.826 to 0.857 (an increase of approximately 3\%) and the accuracy from 0.592 to 0.669 (an increase of about 7\%) compared to the \textit{baseline}. 

In comparing architectures within the EfficientNet family, we observed that increasing model complexity did not improve performance, with a slight degradation even in the deeper variants. These results indicate that more complex architectures require more refined training adjustments to reach their full potential. Despite the strong results, the main limitation of this study is the use of a single, albeit multicenter and heterogeneous, dataset, highlighting the need for future validation on external datasets to confirm the robustness and generalizability of the proposed method.

As future work, we propose investigating more recent architectures based on \textit{Transformers}, particularly \textit{Vision Transformers} (ViT) and hybrid \textit{CNN-Transformer} models, which may better capture global dependencies in histopathological images. Additionally, \textit{ensemble learning} techniques may be explored to improve the model's robustness and generalization. These approaches have the potential to further improve system performance, particularly in classifying intermediate ISUP grades, thereby expanding their applicability as diagnostic decision-support tools across diverse clinical contexts.